\documentclass[a4paper,11pt,oneside]{memoir}

% Castellano
\usepackage[spanish,es-tabla]{babel}
\selectlanguage{spanish}
\usepackage[utf8]{inputenc}
\usepackage{placeins}

\RequirePackage{booktabs}
\RequirePackage[table]{xcolor}
\RequirePackage{xtab}
\RequirePackage{multirow}

% Links
\usepackage[colorlinks]{hyperref}
\hypersetup{
	allcolors = {red}
}

% Ecuaciones
\usepackage{amsmath}

% Rutas de fichero / paquete
\newcommand{\ruta}[1]{{\sffamily #1}}

% Párrafos
\nonzeroparskip

% inglés
\newcommand{\english}[1]{\emph{#1}}

% Símbolo euro
\usepackage{eurosym}
\usepackage[utf8]{inputenc}
\DeclareUnicodeCharacter{20AC}{\euro}

% Float barrier
\usepackage{placeins}
% float
\usepackage{float}

% Algoritmos
\usepackage[spanish]{algorithm2e}

% Footnotes
\usepackage{footnote}

% Imagenes
\usepackage{graphicx}
\newcommand{\imagen}[2]{
	\begin{figure}[!h]
		\centering
		\includegraphics[width=0.9\textwidth]{#1}
		\caption{#2}\label{fig:#1}
	\end{figure}
	\FloatBarrier
}

\newcommand{\imagenflotante}[2]{
	\begin{figure}%[!h]
		\centering
		\includegraphics[width=0.9\textwidth]{#1}
		\caption{#2}\label{fig:#1}
	\end{figure}
}



% El comando \figura nos permite insertar figuras comodamente, y utilizando
% siempre el mismo formato. Los parametros son:
% 1 -> Porcentaje del ancho de página que ocupará la figura (de 0 a 1)
% 2 --> Fichero de la imagen
% 3 --> Texto a pie de imagen
% 4 --> Etiqueta (label) para referencias
% 5 --> Opciones que queramos pasarle al \includegraphics
% 6 --> Opciones de posicionamiento a pasarle a \begin{figure}
\newcommand{\figuraConPosicion}[6]{%
  \setlength{\anchoFloat}{#1\textwidth}%
  \addtolength{\anchoFloat}{-4\fboxsep}%
  \setlength{\anchoFigura}{\anchoFloat}%
  \begin{figure}[#6]
    \begin{center}%
      \Ovalbox{%
        \begin{minipage}{\anchoFloat}%
          \begin{center}%
            \includegraphics[width=\anchoFigura,#5]{#2}%
            \caption{#3}%
            \label{#4}%
          \end{center}%
        \end{minipage}
      }%
    \end{center}%
  \end{figure}%
}

%
% Comando para incluir imágenes en formato apaisado (sin marco).
\newcommand{\figuraApaisadaSinMarco}[5]{%
  \begin{figure}%
    \begin{center}%
    \includegraphics[angle=90,height=#1\textheight,#5]{#2}%
    \caption{#3}%
    \label{#4}%
    \end{center}%
  \end{figure}%
}
% Para las tablas
\newcommand{\otoprule}{\midrule [\heavyrulewidth]}
%
% Nuevo comando para tablas pequeñas (menos de una página).
\newcommand{\tablaSmall}[5]{%
 \begin{table}
  \begin{center}
   \rowcolors {2}{gray!35}{}
   \begin{tabular}{#2}
    \toprule
    #4
    \otoprule
    #5
    \bottomrule
   \end{tabular}
   \caption{#1}
   \label{tabla:#3}
  \end{center}
 \end{table}
}

%
% Nuevo comando para tablas pequeñas (menos de una página).
\newcommand{\tablaSmallSinColores}[5]{%
 \begin{table}[H]
  \begin{center}
   \begin{tabular}{#2}
    \toprule
    #4
    \otoprule
    #5
    \bottomrule
   \end{tabular}
   \caption{#1}
   \label{tabla:#3}
  \end{center}
 \end{table}
}

\newcommand{\tablaApaisadaSmall}[5]{%
\begin{landscape}
  \begin{table}
   \begin{center}
    \rowcolors {2}{gray!35}{}
    \begin{tabular}{#2}
     \toprule
     #4
     \otoprule
     #5
     \bottomrule
    \end{tabular}
    \caption{#1}
    \label{tabla:#3}
   \end{center}
  \end{table}
\end{landscape}
}

%
% Nuevo comando para tablas grandes con cabecera y filas alternas coloreadas en gris.
\newcommand{\tabla}[6]{%
  \begin{center}
    \tablefirsthead{
      \toprule
      #5
      \otoprule
    }
    \tablehead{
      \multicolumn{#3}{l}{\small\sl continúa desde la página anterior}\\
      \toprule
      #5
      \otoprule
    }
    \tabletail{
      \hline
      \multicolumn{#3}{r}{\small\sl continúa en la página siguiente}\\
    }
    \tablelasttail{
      \hline
    }
    \bottomcaption{#1}
    \rowcolors {2}{gray!35}{}
    \begin{xtabular}{#2}
      #6
      \bottomrule
    \end{xtabular}
    \label{tabla:#4}
  \end{center}
}

%
% Nuevo comando para tablas grandes con cabecera.
\newcommand{\tablaSinColores}[6]{%
  \begin{center}
    \tablefirsthead{
      \toprule
      #5
      \otoprule
    }
    \tablehead{
      \multicolumn{#3}{l}{\small\sl continúa desde la página anterior}\\
      \toprule
      #5
      \otoprule
    }
    \tabletail{
      \hline
      \multicolumn{#3}{r}{\small\sl continúa en la página siguiente}\\
    }
    \tablelasttail{
      \hline
    }
    \bottomcaption{#1}
    \begin{xtabular}{#2}
      #6
      \bottomrule
    \end{xtabular}
    \label{tabla:#4}
  \end{center}
}

%
% Nuevo comando para tablas grandes sin cabecera.
\newcommand{\tablaSinCabecera}[5]{%
  \begin{center}
    \tablefirsthead{
      \toprule
    }
    \tablehead{
      \multicolumn{#3}{l}{\small\sl continúa desde la página anterior}\\
      \hline
    }
    \tabletail{
      \hline
      \multicolumn{#3}{r}{\small\sl continúa en la página siguiente}\\
    }
    \tablelasttail{
      \hline
    }
    \bottomcaption{#1}
  \begin{xtabular}{#2}
    #5
   \bottomrule
  \end{xtabular}
  \label{tabla:#4}
  \end{center}
}



\definecolor{cgoLight}{HTML}{EEEEEE}
\definecolor{cgoExtralight}{HTML}{FFFFFF}

%
% Nuevo comando para tablas grandes sin cabecera.
\newcommand{\tablaSinCabeceraConBandas}[5]{%
  \begin{center}
    \tablefirsthead{
      \toprule
    }
    \tablehead{
      \multicolumn{#3}{l}{\small\sl continúa desde la página anterior}\\
      \hline
    }
    \tabletail{
      \hline
      \multicolumn{#3}{r}{\small\sl continúa en la página siguiente}\\
    }
    \tablelasttail{
      \hline
    }
    \bottomcaption{#1}
    \rowcolors[]{1}{cgoExtralight}{cgoLight}

  \begin{xtabular}{#2}
    #5
   \bottomrule
  \end{xtabular}
  \label{tabla:#4}
  \end{center}
}




\graphicspath{ {./img/} }

% Capítulos
\chapterstyle{bianchi}
\newcommand{\capitulo}[2]{
	\setcounter{chapter}{#1}
	\setcounter{section}{0}
	\chapter*{#2}
	\addcontentsline{toc}{chapter}{#2}
	\markboth{#2}{#2}
}

% Apéndices
\renewcommand{\appendixname}{Apéndice}
\renewcommand*\cftappendixname{\appendixname}

\newcommand{\apendice}[1]{
	%\renewcommand{\thechapter}{A}
	\chapter{#1}
}

\renewcommand*\cftappendixname{\appendixname\ }

% Formato de portada
\makeatletter
\usepackage{xcolor}
\newcommand{\tutor}[1]{\def\@tutor{#1}}
\newcommand{\course}[1]{\def\@course{#1}}
\definecolor{cpardoBox}{HTML}{E6E6FF}
\def\maketitle{
  \null
  \thispagestyle{empty}
  % Cabecera ----------------
\noindent\includegraphics[width=\textwidth]{cabecera}\vspace{1cm}%
  \vfill
  % Título proyecto y escudo informática ----------------
  \colorbox{cpardoBox}{%
    \begin{minipage}{.8\textwidth}
      \vspace{.5cm}\Large
      \begin{center}
      \textbf{TFG del Grado en Ingeniería Informática}\vspace{.6cm}\\
      \textbf{\LARGE\@title{}}
      \end{center}
      \vspace{.2cm}
    \end{minipage}

  }%
  \hfill\begin{minipage}{.20\textwidth}
    \includegraphics[width=\textwidth]{escudoInfor}
  \end{minipage}
  \vfill
  % Datos de alumno, curso y tutores ------------------
  \begin{center}%
  {%
    \noindent\LARGE
    Presentado por \@author{}\\ 
    en Universidad de Burgos --- \@date{}\\
    Tutores: \@tutor{}\\
  }%
  \end{center}%
  \null
  \cleardoublepage
  }
\makeatother


% Datos de portada
\title{Extracción y procesamiento de datos de Amazon para su utilización en un estudio de marketing.}
\author{Daniel Arnaiz Gutierrez}
\tutor{Dr. José Francisco Díez Pastor \\Dr. César Ignacio García Osorio}
\date{\today}

\begin{document}

\maketitle



\cleardoublepage



%%%%%%%%%%%%%%%%%%%%%%%%%%%%%%%%%%%%%%%%%%%%%%%%%%%%%%%%%%%%%%%%%%%%%%%%%%%%%%%%%%%%%%%%



\frontmatter


\clearpage

% Indices
\tableofcontents

\clearpage

\listoffigures

\clearpage

\listoftables

\clearpage

\mainmatter

\appendix

\apendice{Plan de Proyecto Software}

\section{Introducción}

La planificación de un proyecto es quizás la fase mas importante de éste. Una buena planificación consigue que el desarrollo del proyecto siga su curso con normalidad y con el mínimo número de imprevistos. A la hora de planificar el proyecto, son varias las partes a tener en cuenta: tiempo, trabajo, y dinero.

A continuación se tratarán los detalles del plan de proyecto llevado a cabo en dos secciones:

\begin{description}
    \item[Planificación temporal:] En este apartado se analizará y explicará cómo se ha planificado el tiempo durante el desarrollo del proyecto teniendo en cuenta el trabajo necesario para cada parte.
    \item[Estudio de viabilidad:] En este apartado se estudiará la viabilidad del proyecto en distintos campos:
    \begin{description}
        \item[Viabilidad económica:] Estimación de los costes y beneficios del proyecto.
        \item[Viabilidad legal:] Análisis sobre los conceptos legales del proyecto, esto engloba las licencias y políticas utilizadas.
    \end{description}
\end{description}

\section{Planificación temporal}

Antes de empezar con el análisis sobre la planificación temporal es necesario añadir que se ha seguido una metodología ágil para la realización del proyecto, más en concreto se ha seguido el método \english{Scrum}. Es por esto que el desarrollo del proyecto se ha dividido en diferentes \english{sprints}, con reuniones entre todos los integrantes del proyecto para tratar los cambios y las siguientes tareas a realizar.

A continuación se mostrarán las principales características de cada \english{sprint} junto con un gráfico \english{burndown} generado con la ayuda de la extensión ZenHub y GitHub.

\subsection{\english{Sprint} 1} %(18/11/2018 - 24/3/2019)}
El primer \english{sprint} ha resultado ser el más largo. Esto se debe a que el desarrollo del proyecto comenzó mas tarde de lo que en un principio se había planeado. Como se puede ver en el gráfico \english{burndown} a continuación, en el último tramo del \english{sprint} es cuando de verdad las tareas propuestas comienzan a darse por completadas.

\imagen{anexos/sprint1}{Gráfico \english{burndown} del \english{sprint} 1}

Este primer \english{sprint} se centra principalmente en la puesta en marcha de las principales herramientas a utilizar, además de la toma de algunas decisiones como pueden ser la elección de qué \english{web scraper} y etiquetador de imágenes utilizar.

Debido a la extensa duración de este \english{sprint}, al finalizar ya encontramos algunas tareas más avanzadas completadas:

\begin{itemize}
    \item Primeras versiones de \english{spiders} funcionales.
    \item Etiquetado manual de un número reducido de imágenes.
    \item Extracción de los primeros productos con sus respectivos campos.
\end{itemize}


\subsection{\english{Sprint} 2} %(17/3/2018 - 12/4/2019)}
 Este \english{sprint} cuenta con una duración mucho menor al anterior. Es aquí donde comienza la familiarización con \LaTeX{} y las primeras mejoras al \english{web scraper}.
 
\imagen{anexos/sprint2}{Gráfico \english{burndown} del \english{sprint} 2}
 
Las tareas realizadas en este \english{sprint} han sido las siguientes:
 
\begin{itemize}
    \item Prueba con nuevas etiquetas en el clasificador manual de imágenes.
    \item Comenzar con la documentación del proyecto.
    \item Modificación de los campos a extraer en el \english{web scraper}.
\end{itemize}


\subsection{\english{Sprint} 3} %(12/4/2018 - 3/5/2019)}
Este \english{sprint} se centra en añadir nuevas funcionalidades al proyecto, como pueden ser el almacenamiento en la base datos o la extracción de comentarios.

\imagen{anexos/sprint3}{Gráfico \english{burndown} del \english{sprint} 3}

Las tareas realizadas en este \english{sprint} han sido las siguientes:
 
\begin{itemize}
    \item Extracción de comentarios asociados a cada producto.
    \item Comenzar a hacer pruebas con el almacenamiento en una base de datos.
    \item Estudiar como generar un documento Excel a partir de los datos extraídos.
    \item Buscar información sobre como aplicar regresión lineal sobre la información extraída.
\end{itemize}

\subsection{\english{Sprint} 4} %(3/5/2018 - 20/5/2019)}
Este \english{sprint} ha estado orientado sobre todo a la investigación de posibles modificaciones y mejoras.

\imagen{anexos/sprint4}{Gráfico \english{burndown} del \english{sprint} 4}

Las tareas realizadas en este \english{sprint} han sido las siguientes:
 
\begin{itemize}
    \item Investigación sobre OpenCV y sus librerías de detección facial.
    \item Primeras iteraciones del clasificador automático de imágenes.
    \item Probar el funcionamiento del detector facial sobre nuestro propio \english{dataset}.
    \item Añadir los enlaces de productos para mujeres en el \english{web scraper}.
    \item Añadir un nuevo campo con el sexo al que va dirigido un artículo a la hora de extraer los productos.
\end{itemize}


\subsection{\english{Sprint} 5} %(22/5/2018 - 15/6/2019)}
\english{Sprint} algo más extenso que se centra principalmente en la primera versión funcional del clasificador automático y la resolución de algunos problemas encontrados.

\imagen{anexos/sprint5}{Gráfico \english{burndown} del \english{sprint} 5}

Las tareas realizadas en este \english{sprint} han sido las siguientes:

\begin{itemize}
    \item Extracción de 1000 productos para la creación del \english{dataset}a final.
    \item Primera versión funcional del clasificador automático Modelo/No modelo.
    \item Primeras pruebas con la detección de caras parcialmente visibles.
    \item Investigar sobre como generar un domuento Excel a partir de la base de datos.
\end{itemize}


\subsection{\english{Sprint} 6} %(16/6/2018 - 30/6/2019)}
\english{Sprint} final del proyecto. Es aquí donde se finaliza el proyecto y se hace casi la totalidad de la documentación.

\imagen{anexos/sprint5}{Gráfico \english{burndown} del \english{sprint} 6}

Las tareas realizadas en este \english{sprint} han sido las siguientes:

\begin{itemize}
    \item Separación del campo que recoge el rango de precios de un producto en dos: precio mínimo y precio máximo.
    \item Revisión y actualización de los nombres de las tablas y los campos de la base de datos.
    \item Finalización de la documentación del proyecto y sus respectivos anexos.
    \item Finalización de los \english{notebooks} de entrenamiento y clasificación de imágenes.
\end{itemize}



\section{Estudio de viabilidad}

\subsection{Viabilidad económica}

En este apartado se analizará los costes y beneficios estimados de este proyecto en un entorno empresarial.

\subsection{Costes}
\subsubsection{Costes de personal}
La totalidad de este proyecto ha sido desarrollada por una única persona a tiempo parcial durante un periodo aproximado de 5 meses. Aplicando el salario mínimo y las cuotas de la Seguridad social\footnote{\url{http://www.seg-social.es/wps/portal/wss/internet/Trabajadores/CotizacionRecaudacionTrabajadores/36537}}, se puede estimar un coste de personal de la siguiente forma:

\begin{table}[!h]
	\centering
	\begin{tabular}{@{}l|l@{}}
		\toprule
		Concepto & Coste (\euro) \\
		\midrule
		Salario neto & \EUR{1000}  \\
		Retención IRPF (19 \%) & \EUR{360,53} \\
		Seguridad social (28,30 \%) & \EUR{537,00} \\
		\midrule
		Salario total (mensual) & \EUR{1897,53} \\
		\midrule
		Total 5 meses & \EUR{9.487,65} \\
		\bottomrule
	\end{tabular}
	\caption{Costes de personal}
	\label{tab:personal}
\end{table}

\subsubsection{Costes de material}

En cuanto a los costes materiales y a nivel de programas informáticos, el único gasto es el del ordenador utilizado para la realización del proyecto. Todas las librerías utilizadas son gratuitas y no es necesario la contratación de un servidor para el funcionamiento del proyecto.

Por lo tanto, asociando un valor aproximado de 1300\euro\ al ordenador y estimando una amortización de 5 años, el coste total de materiales para estos 5 meses es el siguiente:

$$\dfrac{\textup{1.300 \euro}}{5\ años * 12\ meses} = \textup{21,67 \euro} $$

\begin{table}[!h]
	\centering
	\begin{tabular}{@{}l|l|l@{}}
		\toprule
		\textbf{Concepto} & \textbf{Coste} & \textbf{Amortización} \\
		\midrule
		Ordenador & 1.300 \euro & 21,67 \euro \\
		\midrule
		\textbf{Total 5 meses} & 108,34 \euro \\
		\bottomrule
	\end{tabular}
	\caption{Costes de hardware}
	\label{tab:hardware}
\end{table}


\subsubsection{Coste total}

A continuación se muestra el coste total del proyecto: 

\FloatBarrier
\begin{table}[h]
	\centering
	\begin{tabular}{@{}l|l@{}}
		\toprule
		\textbf{Concepto} & \textbf{Coste} \\
		\midrule
		Personal & 9.487,65 \euro \\
		Material & 108,34 \euro \\
		\midrule
		\textbf{Total} & 9.595,99 \euro \\
		\bottomrule
	\end{tabular}
	\caption{Costes totales del proyecto}
	\label{tab:total}
\end{table}
\FloatBarrier

\subsection{Beneficios}
En cuanto a los beneficios del proyecto, no se desarrollado teniendo como meta la generación de ingresos. Se podría comercializar cobrando una tarifa mensual a los clientes u ofreciendo un servicio personalizado a cada cliente que dependería del número de productos a extraer.


\subsection{Viabilidad legal}

Para el estudio de la viabilidad del producto, se van a analizar las librerías usadas en el proyecto, anotando las licencias de las que hacen uso. A continuación se listan:

\begin{savenotes}
\begin{table}[!h]
	\centering
	\begin{tabular}{l|l|l}
		\hline
		Dependencia & Versión & Licencia \\ \hline
		Scrapy & 1.5.1 & BSD \\
		Pillow & 6.0.0 & PIL\footnote{\url{http://www.pythonware.com/products/pil/license.htm}} \\
		requests & 2.21.0 & MIT \\
		js2xml & 0.3.1 & MIT \\
		Keras & 2.3.1 & MIT \\
		Numpy & 1.16.4 & BSD \\
		OpenCV & 4.1.0 & BSD \\
		JSON & \ \ \ \ \ - & FREE \\
		SQLite & 3.28.0 & FREE \\
		Pandas & 0.24.2 & BSD \\
		\bottomrule
	\end{tabular}
	\caption{Licencias utilizadas}
	\label{tab:licencias}
\end{table}
\end{savenotes}
\apendice{Especificación de Requisitos}

\section{Introducción}

\section{Objetivos generales}

\section{Catalogo de requisitos}

\section{Especificación de requisitos}



\apendice{Especificación de diseño}

\section{Introducción}

Este apartado se encarga de recoger los diferentes diseños que han sido llevados a cabo para la realización del proyecto y cumplir de forma satisfactoria los requisitos y objetivos anteriormente tratados. 

\section{Diseño de datos}

Esta sección detalla cómo se han almacenado los datos, tanto en la base de datos SQLite, como en los archivos JSON que genera Scrapy.

\subsection{Estructura JSON}
Cada vez que \english{scrapeamos} un producto, éste será añadido a un archivo JSON que genera Scrapy. En total, se generan 3 archivos JSON por cada uso de la herramienta:

\begin{itemize}
    \item \textbf{Enlaces a cada producto:} El primer JSON recoge todos los enlaces que apuntan a los productos de los cuales se desea extraer la información. La estructura es muy sencilla ya que únicamente guarda un enlace por cada producto:
    
    \imagen{anexos/json_url}{Estructura del JSON con los enlaces a los productos}
    
    Como se puede ver, cada página de <<Amazon.com>> contiene 48 productos. En este JSON se almacenan los enlaces a estos 48 artículos.
    
    \item \textbf{Campos principales de los productos:} Este JSON almacena los campos que se extraen de cada producto de la siguiente forma:
    
    \imagen{anexos/json_product}{Estructura del JSON con los campos de cada producto}
    
    Como se puede apreciar, la estructura se compone de los siguientes campos:
    
    \begin{description}
        \item[ASIN:] Código identificativo de cada producto.
        \item[Sexo:] Sexo al que va dirigido cada artículo.
        \item[Rango de precios:] El precio de cada producto puede variar en función del color o la talla. Este campo recoge el precio máximo y mínimo de cada producto.
        \item[Puntuación:] Media de las puntuaciones que los clientes han votado a cada producto. Va del 0 al 5.
        \item[Número de valoraciones:] Número de clientes que han proporcionado una valoración del producto.
        \item[Marca:] El nombre de la marca de cada artículo.
        \item[Descripción:] Texto que el vendedor ha utilizado para describir y dar detalles de cada producto.
        \item[Enlaces a las imágenes:] Enlaces de las imágenes asociadas a cada artículo.
    \end{description}
    
    \item \textbf{Comentarios de los productos:} Por último, este JSON se encarga de recoger los diferentes comentarios que los clientes han dejado a cada artículo. Se compone de dos únicos campos: una lista con los comentarios y el identificador del producto:
    
    \imagen{anexos/json_comment}{Estructura del JSON con los comentarios de cada producto}
\end{itemize}

\subsection{Estructura de la base de datos SQLite}

Este proyecto también hace uso de una base de datos para almacenar toda la información una vez han sido procesadas y clasificadas las imágenes.
Para ello se hace uso de una base de datos cuyas tablas y campos son los siguientes:

\imagen{bd_diagram}{Diagrama de la base de datos utilizada.}

\begin{description}
    \item[Tabla productos] Recoge los principales campos de cada producto. Cada entrada pertenece a un artículo diferente.
    \item[Tabla Imágenes] Almacena los enlaces a las imágenes de cada producto. Puede haber varias imágenes para un único producto.
    \item[Tabla Comentarios] Almacena los comentarios asociados a cada artículo. Puede haber varios comentarios de cada producto.
    \item[Tabla Predicciones] Almacena las predicciones que los clasificadores han generado para la imagen principal de cada artículo. Cada entrada es de un producto diferente.
\end{description}

\section{Diseño procedimental}
En esta sección se muestran los diferentes procedimientos para la realización de los procesos mas importantes del proyecto:


\subsection{\english{Web scraping}}

A la hora de la extracción de información de los productos, el proceso a seguir es el siguiente:

\begin{enumerate}
    \item Extraer los enlaces a dichos productos. Para ello se hace uso del \english{spider} llamado <<urls\_spider>> en Scrapy.
    \item Con el archivo JSON que se ha generado en el anterior paso, se pasará a la extracción de información de cada producto por medio del \english{spider} llamado <<products\_spider>>.
    \item Por último, si se desean extraer los comentarios de los productos, se hará uso del \english{spider} llamado <<comments\_spider>> para ello.
\end{enumerate}

\imagen{proyecto_scrapy}{Proceso de \english{web scraping} del proyecto}

\subsection{Entrenamiento de los clasificadores}

Para entrenar los clasificadores es necesario haber creado un conjunto de entrenamiento y haber etiquetado de forma manual las imágenes que conforman dicho conjunto.
Una vez se dispone de ese conjunto de entrenamiento los pasos para entrenar el clasificador son los siguientes:

\begin{enumerate}
    \item Descargar el \english{dataset} que contiene las imágenes previamente etiquetadas.
    \item De forma secuencial, se convertirá cada imagen a formato \english{array} y se almacenará junto a su etiqueta.
    \item Crear conjuntos de test y entrenamiento para aplicar la técnica de validación cruzada.
    \item Inicializar el modelo del clasificador.
    \item Entrenar el clasificador.
    \item Almacenar el clasificador ya entrenado para su posterior uso.
\end{enumerate}

\subsection{Clasificación de imágenes y almacenamiento de resultados}

El proceso de clasificación de imágenes y su posterior almacenamiento es el siguiente:

\begin{enumerate}
    \item Importar y cargar los clasificadores previamente entrenados.
    \item Importar el archivo JSON que contiene los productos a clasificar.
    \item Descargar cada imagen con el ASIN del producto como nombre.
    \item Clasificar cada imagen siguiendo el siguiente algoritmo:
    
        \FloatBarrier
        \begin{algorithm}[!h]
            \If{Imagen cargada correctamente}
            	{
            	Redimensionar imagen.
            	
            	Convertir imagen a formato \english{array}.
            	
            	Aplicar clasificador que detecta modelo.
            	
            	\If{El clasificador ha encontrado un modelo en la imagen}
            		{
            			Aplicar clasificador que detecta cara.
            			
            			\If{El clasificador ha encontrado la cara del modelo}
            		    {
            			Devolver predicción: Modelo con cara.
            		    }
            		    Devolver predicción: Modelo sin cara.
            		}
            	    Devolver predicción: Sin modelo.
            	}
            \caption{Algoritmo de clasificación de imágenes}
            \label{Alg:}
        \end{algorithm}
        \FloatBarrier
        
    \item Importar la base de datos
    \item Crear la tabla <<Predicciones>> en la base de datos
    \item Almacenar en la base de datos cada predicción generada por el clasificador quedando asociada al artículo al que pertenece la imagen.
            
\end{enumerate}

\section{Diseño arquitectónico}
 En cuanto al diseño arquitectónico del proyecto, es necesario señalar que la primera parte del proyecto, todo lo relacionado con la extracción de datos y \english{web scraping}, funciona bajo la estructura se Scrapy. El resto se encuentra dividido en dos \english{notebooks} de Google Colab, el primero para el entrenamiento de los clasificadores, y el segundo destinado a la clasificación de imágenes y su posterior almacenamiento en la base de datos.
 
El siguiente diagrama muestra una visión general de la arquitectura de Scrapy junto con los conponentes que lo conforman y el flujo de datos que se lleva a cabo en el sistema \cite{scrapy:arch}:
 
 \imagen{scrapy_architecture}{Arquitectura que sigue Scrapy y sus diferentes componentes}
 
 \subsection{Componentes de scrapy}
 
 A continuación se describen los principales componentes de la arquitectura de un proyecto de Scrapy:
 
\begin{itemize}
    \item \textbf{\english{Items}:} Los diferentes elementos o campos a extraer.
    \item \textbf{\english{Spiders}:} Clases escritas por el usuario que contienen los diferentes procedimientos para parsear y extraer los \english{items} de uno o varios enlaces. En este caso hemos hecho uso de 3 \english{spiders}: Una para obtener los enlaces a cada artículo, otra para extraer los comentarios, y la última para el resto de campos.
    \item \textbf{\english{Pipelines}:} Apartado responsable de procesar los \english{items} una vez han sido extraidos por las \english{spiders}. En este proyecto se han usado para almacenar los campos extraídos en una base de datos.
\end{itemize}

%\imagen{anexos/arq}{Estructura del proyecto de scrapy}
\apendice{Documentación técnica de programación}

\section{Introducción}

Esta sección se encarga de detallar los conceptos relacionados con la documentación técnica que sigue el proyecto. Se tratará de explicar la estructura de directorios que sigue el proyecto, así como los pasos a seguir para ser instalado y ejecutado.

\section{Estructura de directorios}

A continuación se muestra la estructura de directorios en la que se distribuye el proyecto:

\begin{itemize}
    \item \textbf{Scrapy/} Directorio principal de Scrapy. Se compone de los siguientes archivos y subdirectorios:
        \begin{itemize}
            \item \textbf{amazon/spiders/} Directorio donde se encuentran las diferentes \english{spiders}.
                \begin{itemize}
                    \item \textbf{urls\_spiders.py} \english{Spider} utilizada para extraer los enlaces de los productos.
                    \item \textbf{products\_spiders.py} \english{Spider} utilizada para extraer los principales campos de cada producto.
                    \item \textbf{comments\_spiders.py} \english{Spider} utilizada para extraer los comentarios asociados a cada producto.
                \end{itemize}
            \item \textbf{amazon/items.py} Fichero que recoge los diferentes campos que extrae Scrapy.
            \item \textbf{amazon/middlewares.py} Fichero que recoge la configuración de los \english{middlewares} se Scrapy si los hubiera.
            \item \textbf{amazon/pipelines.py} Fichero que se encarga del almacenamiento en la base de datos de los campos extraídos.
            \item \textbf{amazon/settings.py} Fichero de configuración del proyecto de Scrapy.
        \end{itemize}
    \item \textbf{Documentacion/} Este directorio contiene todos los archivos relacionados con la documentación y la memoria del proyecto, Tanto los ficheros fuente de \LaTeX{}, como los PDFs ya finalizados.
    \item \textbf{Colab notebooks/} Directorio donde se encuentran los \english{notebooks} de Google Colab utilizados en el proyecto:
        \begin{itemize}
            \item \textbf{Clasificador\_Modelo\_Sin\_modelo.ipynb} \english{Notebook} encargado del entrenamiento del clasificador de  imágenes que detecta si hay un modelo en la imágen o no.
            \item \textbf{Clasificador\_cara\_sin\_cara.ipynb} \english{Notebook} encargado del entrenamiento del clasificador de  imágenes que detecta si la cara de un modelo es visible o no.
            \item \textbf{Clasificador\_final\_excel\_db.ipynb} \english{Notebook} encargado de la clasificación de imágenes y su posterior almacenamiento en la base de datos y documento Excel.
        \end{itemize}
    \item \textbf{Databases/} Carpeta que contiene las diferentes bases de datos utilizadas durante la realización del proyecto.
    \item \textbf{Datasets/} Directorio que recoge los \english{datasets} utilizados para el entrenamiento de los clasificadores:
        \begin{itemize}
            \item \textbf{dataset-modelo.zip} \english{Dataset} que contiene las imágenes ya etiquetadas utilizadas para el entrenamiento del clasificador de imágenes que detecta modelos.
            \item \textbf{dataset-cara.zip} \english{Dataset} que contiene las imágenes ya etiquetadas utilizadas para el entrenamiento del clasificador de imágenes que detecta caras.
        \end{itemize}
    \item \textbf{Keras models/} Directorio donde encontramos los clasificadores de Keras ya entrenados para su posterior utilización:
        \begin{itemize}
            \item \textbf{clasificador-modelo.zip} Clasificador de imágenes encargado de detectar modelos ya entrenado y listo para su utilización.
            \item \textbf{clasificador-cara.zip} Clasificador de imágenes encargado de detectar caras ya entrenado y listo para su utilización.
        \end{itemize}
    \item \textbf{Outputs/} Directorio donde se guardan todos los resultados del proyecto: Archivos JSON y documentos Excel.
    \item \textbf{Scripts/} Directorio que incluye algunos \english{scripts} utilizados para algunas tareas del proyecto:
        \begin{itemize}
            \item \textbf{JSONmerger.py} \english{Script} que combina los campos de dos archivos JSON utilizando el campo <<url>> como elemento identificador.
            \item \textbf{dataextractor.py} \english{Script} que descarga imágenes dada su URL y las renombra en funcion de cómo han sido etiquetadas. Utilizado para la creación de los \english{datasets}.
            \item \textbf{dataturks\_fixer.py} \english{Script} usado para arreglar el JSON generado por Dataturks para su correcto funcionamiento.
            \item \textbf{url\-extractor.py} \english{Script} que extrae los enlaces de las imágenes de cada producto y lo añade a una lista para posteriormente etiquetarlo en Dataturks.
        \end{itemize}
\end{itemize}

\section{Manual del programador}

Apartado que detalla algunos de los procesos llevados a cabo en el proyecto necesarios para que pueda seguir desarrollándose.

\subsection{Manual de Scrapy}

Para la creación del proyecto de Scrapy es necesaria la instalación de la librería. Esto se puede hacer de varias formas, tal y como describe la documentación oficial\footnote{\url{http://doc.scrapy.org/en/latest/intro/install.html}}:

\begin{itemize}
    \item \textbf{Utilizando Conda:} Se instala la librería con el siguiente comando:
        \begin{verbatim}
        conda install -c conda-forge scrapy
        \end{verbatim}
    \item \textbf{Utilizando PIP:} Se instala la librería con el siguiente comando:
        \begin{verbatim}
        pip install Scrapy
        \end{verbatim}
\end{itemize}

Una vez instalado el paquete de Scrapy, la creación de un nuevo proyecto es sencilla:

\begin{verbatim}
    scrapy startproject [nombre_proyecto]
\end{verbatim}

Este comando se encarga de crear la estructura de carpetas necesaria para el correcto funcionamiento de Scrapy. A continuación se detalla el proceso a seguir para añadir \english{spiders} al proyecto:

\begin{verbatim}
    cd [nombre_proyecto]
    scrapy genspider [nombre_spider] [direccion_web]
\end{verbatim}

En cuanto a la ejecución de un \english{spider}, se hace con el siguiente comando:

\begin{verbatim}
    scrapy crawl [nombre_spider] -o [nombre_output]
\end{verbatim}

\subsection{Manual de Dataturks}

A continuación se detallará el proceso a seguir para el etiquetado manual de imágenes con la ayuda de la herramienta Dataturks\footnote{\url{https://dataturks.com/}}.

\begin{enumerate}
    \item Iniciar sesión por medio de GiHub o creando una cuenta a parte.
    \item Crear un nuevo \english{dataset} y elegir el tipo adecuado, en este caso será el llamado \english{Image Classification}.
    
    \imagen{anexos/dataturks_proyecto}{Tipo de \english{dataset} a crear en Dataturks}
    
    \item Configurar el funcionamiento del \english{dataset}
    
    \imagen{anexos/dataturks_inst}{Creación de un \english{dataset} en Dataturks}
    
    \item Importar una lista con los enlaces a las imágenes que se desean etiquetar. Deberán estar en un documento \english{txt}.
    
    \item Una vez subidas las imágenes, se puede empezar a etiquetar seleccionando la etiqueta adecuada en cada caso.
    
    \imagen{dataturks}{Proceso de etiquetado manual en Dataturks}
    
    \item Por último, para exportar los resultados simplemente se ha de descargar el archivo generado pulsando en el botón \english{Download}.
    
\end{enumerate}

\subsection{Entrenamiento de los clasificadores}

El proceso de entrenamiento de los clasificadores es muy sencillo ya que está muy automatizado. Los \english{notebooks} se encargan de descargar y descomprimir el \english{dataset} previamente subido al repositorio del proyecto para así entrenar el clasificador. Una vez entrenado, se muestra la precisión alcanzada y un resumen de las capas utilizadas por el modelo de Keras.

Por último se descarga el clasificador ya entrenado para su posterior uso.

\section{Compilación, instalación y ejecución del proyecto}

Lo más importante a la hora de instalar y ejecutar este proyecto es la parte de Scrapy. Puesto que en el repositorio ya se encuentra el proyecto de Scrapy al completo, lo primero será descargar el contenido del repositorio. Para ello se utiliza el siguiente comando:

\begin{verbatim}
    git clone https://github.com/daniarnaizg/TFG-Amazon-Scraper.git
\end{verbatim}

Con esto, quedará descargado el repositorio. Lo siguiente será instalar el paquete de Scrapy si no se tiene instalado (Ver manual de Scrapy en la sección anterior).

También es necesaria la instalación de dos paquetes, se pueden instalar por medio de PIP:

\begin{verbatim}
    pip install js2xml
    pip install pypiwin32
\end{verbatim}

A partir de aquí simplemente se puede lanzar una \english{spider}, tal y cómo se ha mostrado antes. Es necesario indicar el fichero al que se desea exportar el contenido descargado en formato JSON:

\begin{verbatim}
    scrapy crawl urls_spider -o urls_output.json
    scrapy crawl products_spider -o products_output.json
    scrapy crawl comments_spider -o comments_output.json
\end{verbatim}

En cuanto a la ejecución de los \english{notebooks}, esto se hará importandolos en el entondo de Google Colab a través del siguiente enlace:
\newline
\newline
\url{https://colab.research.google.com/}


%\section{Pruebas del sistema}

\apendice{Documentación de usuario}

\section{Introducción}

\section{Requisitos de usuarios}

\section{Instalación}

\section{Manual del usuario}





\bibliographystyle{plain}
\bibliography{bibliografiaAnexos}

\end{document}
